\documentclass[12pt, a4paper]{article}

% --- 1. Cấu hình Tiếng Việt và Font chữ ---
\usepackage[utf8]{inputenc}
\usepackage[T5]{fontenc}
\usepackage[vietnamese]{babel}
\usepackage{newtxtext,newtxmath}

\makeatletter
\renewcommand{\normalsize}{\@setfontsize\normalsize{13pt}{16pt}}
\makeatother
\normalsize

% 2. Gói Toán học & Ký hiệu bổ trợ
\usepackage{amsmath, amssymb, amsfonts, mathrsfs, amsthm}
\usepackage{graphicx}
\usepackage{fontawesome}
\usepackage{booktabs, tabularx, array, multirow}
\usepackage{enumitem}

% 3. Đồ họa TikZ, Bảng biến thiên & Hình không gian
\usepackage{tikz}
\usetikzlibrary{calc, intersections, angles, quotes, patterns, arrows.meta, 3d}
\usepackage{pgfplots}
\pgfplotsset{compat=1.18}
\usepackage{tkz-euclide}
\usepackage{tkz-tab}

\renewcommand{\vec}[1]{\overrightarrow{#1}}

% 4. Hộp sư phạm (Tcolorbox)
\usepackage[many]{tcolorbox}
\tcbuselibrary{skins, breakable}

\newtcolorbox{pedabox}[1][]{
    enhanced,
    colback=blue!5!white,
    colframe=blue!75!black,
    fonttitle=\bfseries\sffamily,
    title=#1,
    boxrule=0.8pt,
    arc=2mm,
    breakable,
    left=3mm, right=3mm, top=2mm, bottom=2mm
}

\newtcolorbox{scaffoldbox}[1][]{
    enhanced,
    colback=orange!5!white,
    colframe=orange!85!black,
    fonttitle=\bfseries\sffamily,
    title=#1,
    boxrule=0.8pt,
    arc=2mm,
    breakable,
    left=3mm, right=3mm, top=2mm, bottom=2mm
}

\newtcolorbox{aibox}[1][]{
    enhanced,
    colback=green!5!white,
    colframe=teal!80!black,
    fonttitle=\bfseries\sffamily,
    title=#1,
    boxrule=0.8pt,
    arc=2mm,
    breakable,
    left=3mm, right=3mm, top=2mm, bottom=2mm
}

% 5. Căn lề chuẩn văn bản giáo án
\usepackage{geometry}
\geometry{
    a4paper,
    left=25mm,
    right=15mm,
    top=20mm,
    bottom=20mm
}

% 6. Header / Footer chuẩn hóa (BẮT BUỘC CHÍNH XÁC NỘI DUNG NÀY)
\usepackage{fancyhdr}
\usepackage{lastpage}
\pagestyle{fancy}
\fancyhf{}

\fancyhead[L]{\small \textit{Kế hoạch bài dạy Toán 11}}
\fancyhead[R]{\textbf{Trang \thepage/\pageref{LastPage}}}
\fancyfoot[L]{\small \textit{Tổ Toán - Tin}}
\fancyfoot[R]{\small \textit{Trường THCS\&THPT Hồng Vân}}
\renewcommand{\headrulewidth}{0.4pt}
\renewcommand{\footrulewidth}{0.4pt}

% 7. Giãn dòng & Giãn đoạn theo chuẩn Nghị định 30/2020/NĐ-CP
\usepackage{setspace}
\setstretch{1.15}
\setlength{\parindent}{1.0cm}
\setlength{\parskip}{5pt}

\begin{document}

\begin{center}
    {\Large \textbf{KẾ HOẠCH BÀI DẠY MÔN TOÁN LỚP 11}}\\[0.4em]
    {\LARGE \textbf{BÀI 14. PHÉP CHIẾU SONG SONG. HÌNH BIỂU DIỄN}}\\[0.3em]
    {\LARGE \textbf{CỦA MỘT HÌNH KHÔNG GIAN}}\\[0.5em]
    \textbf{Môn học:} Toán 11 | \textbf{Thời lượng:} 02 tiết (Tiết 37, 38 theo PPCT)
\end{center}

\vspace{0.5em}
\hrule
\vspace{1em}

\section*{I. MỤC TIÊU DẠY HỌC}

\subsection*{1. Kiến thức, Năng lực}

\textbf{a) Năng lực Toán học đặc thù (nguyên văn chuẩn CT GDPT 2018):}
\begin{itemize}[leftmargin=1.5em]
    \item \textbf{Năng lực tư duy và lập luận toán học:}
    \begin{itemize}
        \item Nhận biết được khái niệm phép chiếu song song: mặt phẳng chiếu, phương chiếu, tia chiếu, hình chiếu song song của một điểm, một đường thẳng, một đoạn thẳng và một hình hình học.
        \item Hiểu rõ và giải thích được các tính chất cơ bản của phép chiếu song song:
        \begin{itemize}
            \item Biến ba điểm thẳng hàng thành ba điểm thẳng hàng và bảo toàn thứ tự giữa các điểm đó.
            \item Biến đường thẳng thành đường thẳng (nếu không song song với phương chiếu), biến tia thành tia, biến đoạn thẳng thành đoạn thẳng.
            \item Biến hai đường thẳng song song thành hai đường thẳng song song hoặc trùng nhau.
            \item Bảo toàn tỉ số độ dài của hai đoạn thẳng cùng nằm trên một đường thẳng hoặc nằm trên hai đường thẳng song song.
        \end{itemize}
        \item Hiểu rõ định nghĩa và quy tắc xác định hình biểu diễn của một hình không gian trên mặt phẳng vẽ.
    \end{itemize}
    \item \textbf{Năng lực giải quyết vấn đề toán học:}
    \begin{itemize}
        \item Xác định được ảnh của một điểm, một đoạn thẳng, một tam giác, một hình bình hành, một đường tròn qua phép chiếu song song.
        \item Vẽ được hình biểu diễn chính xác của các hình phẳng quen thuộc (tam giác, hình thang, hình bình hành, hình chữ nhật, hình vuông, hình thoi, hình tròn) và các hình khối cơ bản (hình tứ diện, hình chóp, hình lăng trụ, hình hộp).
        \item Tuân thủ nghiêm ngặt quy tắc nét thấy (vẽ nét liền) và nét khuất (vẽ nét đứt) khi vẽ hình không gian.
    \end{itemize}
    \item \textbf{Năng lực mô hình hóa toán học:}
    \begin{itemize}
        \item Sử dụng được kiến thức về phép chiếu song song để giải thích các hiện tượng quang học thực tế: bóng của cây cối, cột cờ, tòa nhà dưới ánh sáng mặt trời (các tia sáng mặt trời xem như song song); ứng dụng trong vẽ phối cảnh kĩ thuật và kiến trúc.
    \end{itemize}
    \item \textbf{Năng lực giao tiếp toán học:}
    \begin{itemize}
        \item Đọc hiểu các bản vẽ hình học không gian, mô tả chính xác mối quan hệ giữa vật thể ba chiều thực tế và hình biểu diễn hai chiều trên mặt phẳng giấy.
    \end{itemize}
    \item \textbf{Năng lực sử dụng công cụ, phương tiện học toán:}
    \begin{itemize}
        \item Sử dụng thành thạo thước kẻ, compa, êke để vẽ hình biểu diễn; khai thác phần mềm GeoGebra 3D và các phần mềm đồ họa máy tính để mô phỏng phép chiếu song song và dựng hình khối 3D.
    \end{itemize}
\end{itemize}

\textbf{b) Năng lực số (theo Thông tư 02/2025/TT-BGDĐT):}
\begin{itemize}[leftmargin=1.5em]
    \item \textbf{Mã 3.1NC1a (Tạo lập và mô phỏng hình khối 3D trong phần mềm đồ họa):} Học sinh sử dụng phần mềm hình học động GeoGebra 3D hoặc phần mềm đồ họa 3D trực quan dựng hình biểu diễn của hình chóp và hình hộp; thay đổi góc nhìn camera chiếu (phép chiếu song song Orthographic vs phép chiếu phối cảnh Perspective) để quan sát sự biến dạng góc và độ dài.
\end{itemize}

\textbf{c) Năng lực AI (theo Quyết định 2422/QĐ-BGDĐT):}
\begin{itemize}[leftmargin=1.5em]
    \item \textbf{Mã 11.C3.2 (Ứng dụng AI trong xử lí đồ họa và tạo lập bóng đổ tự động):} Học sinh tìm hiểu cách các mô hình AI tạo sinh hình ảnh (Generative AI) và thuật toán dò tia thời gian thực (Ray-Tracing AI / Neural Rendering) áp dụng nguyên lí của phép chiếu song song từ nguồn sáng vô cực (mặt trời) để tính toán hướng bóng đổ, vùng sáng tối chân thực trên bề mặt vật thể trong game 3D và phim hoạt hình.
\end{itemize}

\textbf{d) Năng lực chung \& Phẩm chất:}
\begin{itemize}[leftmargin=1.5em]
    \item \textbf{Năng lực chung:} Tự chủ và tự học (tự rèn luyện kĩ năng vẽ hình không gian); giao tiếp và hợp tác nhóm khi cùng thảo luận nhận diện hình biểu diễn.
    \item \textbf{Phẩm chất:} Cẩn thận, tỉ mỉ trong từng nét vẽ; trung thực và có trách nhiệm với sản phẩm học tập.
\end{itemize}

\section*{II. THIẾT BỊ DẠY HỌC VÀ HỌC LIỆU}
\begin{itemize}[leftmargin=1.5em]
    \item \textbf{Giáo viên:} Kế hoạch bài dạy chuẩn CV 5512; bài giảng điện tử PowerPoint trực quan; mô hình khung không gian (khung tứ diện, hình chóp, hình lập phương); đèn chiếu tạo bóng (đèn pin/máy chiếu); phần mềm GeoGebra 3D; Phiếu học tập số 1, 2; Rubric đánh giá.
    \item \textbf{Học sinh:} SGK Toán 11, vở ghi; thước kẻ thẳng, êke, compa, bút chì, tẩy; máy tính hoặc điện thoại có kết nối Internet để thực hành mô phỏng số.
\end{itemize}

\section*{III. TIẾN TRÌNH DẠY HỌC (TRỌN VẸN 02 TIẾT)}

\begin{center}
    \framebox[1.0\textwidth]{\parbox{0.96\textwidth}{
        \textbf{CẤU TRÚC PHÂN PHỐI 02 TIẾT DẠY HỌC (TIẾT 37, 38 THEO PPCT):}\\[0.3em]
        $\bullet$ \textbf{Tiết 1 (Tiết 37 PPCT):} Khái niệm phép chiếu song song --- Các tính chất cơ bản của phép chiếu song song --- Năng lực mô hình hóa bóng nắng thực tế.\\[0.3em]
        $\bullet$ \textbf{Tiết 2 (Tiết 38 PPCT):} Hình biểu diễn của một hình không gian --- Quy tắc vàng vẽ hình biểu diễn phẳng --- Năng lực số GeoGebra 3D (Mã 3.1NC1a) --- Năng lực AI dựng bóng đổ Ray-Tracing (Mã 11.C3.2) --- Luyện tập tổng hợp.
    }}
\end{center}

\vspace{0.8em}
\hrule
\vspace{0.8em}

{\large \textbf{TIẾT 1 (TIẾT 37 THEO PHÂN PHỐI CHƯƠNG TRÌNH)}}\\[0.3em]
\textbf{Nội dung:} Khái niệm phép chiếu song song --- Các tính chất cơ bản của phép chiếu song song.

\subsubsection*{1. Hoạt động 1: Mở đầu / Khởi động (05 phút)}
\textbf{Chủ đề:} \textit{Bóng của cột cờ trên sân trường dưới ánh nắng mặt trời}
\begin{itemize}
    \item \textbf{a) Mục tiêu:} Giúp học sinh liên hệ hiện tượng tạo bóng tự nhiên dưới ánh sáng mặt trời với mô hình toán học phép chiếu song song.
    \item \textbf{b) Nội dung:} Vào một buổi sáng nắng đẹp, bóng của cột cờ in trên sân trường là một đoạn thẳng.
    \begin{itemize}
        \item Vì sao các tia nắng mặt trời chiếu xuống mặt đất được coi là song song với nhau?
        \item Mặt sân trường đóng vai trò là gì? Hướng của tia nắng đóng vai trò là gì?
    \end{itemize}
    \item \textbf{c) Sản phẩm:} Học sinh trả lời: Mặt trời ở rất xa nên chùm tia sáng tới mặt đất song song với nhau. Mặt sân là mặt phẳng hứng bóng (mặt phẳng chiếu), hướng tia nắng là phương chiếu. Bóng cột cờ là hình chiếu của cột cờ.
    \item \textbf{d) Tổ chức thực hiện:} GV đặt câu hỏi $\to$ HS quan sát, thảo luận nhanh $\to$ GV dẫn dắt vào bài học.
\end{itemize}

\subsubsection*{2. Hoạt động 2: Hình thành kiến thức mới (25 phút)}

\paragraph{Mục 1: Khái niệm Phép chiếu song song (10 phút)}
\begin{itemize}
    \item Cho mặt phẳng $(\alpha)$ và một đường thẳng $\Delta$ cắt $(\alpha)$.
    \item Với mỗi điểm $M$ trong không gian:
    \begin{itemize}
        \item Qua $M$ vẽ đường thẳng $d$ song song với $\Delta$ (hoặc trùng với $\Delta$).
        \item Đường thẳng $d$ cắt mặt phẳng $(\alpha)$ tại điểm $M'$.
        \item Điểm $M'$ được gọi là \textbf{hình chiếu song song} của điểm $M$ trên mặt phẳng $(\alpha)$ theo phương $\Delta$.
    \end{itemize}
    \item Phép đặt tương ứng mỗi điểm $M$ với hình chiếu $M'$ của nó được gọi là \textbf{phép chiếu song song} lên mặt phẳng $(\alpha)$ theo phương $\Delta$.
    \item Mặt phẳng $(\alpha)$ gọi là \textbf{mặt phẳng chiếu}, đường thẳng $\Delta$ quy định \textbf{phương chiếu}.
    \item Cho hình $\mathscr{H}$. Tập hợp các điểm $M'$ là hình chiếu của tất cả các điểm $M \in \mathscr{H}$ được gọi là \textbf{hình chiếu song song} của hình $\mathscr{H}$ trên mặt phẳng $(\alpha)$ (kí hiệu $\mathscr{H}'$).
\end{itemize}

\begin{center}
\begin{tikzpicture}[scale=0.85]
    % Mặt phẳng chiếu (alpha)
    \draw[thick] (-0.5,-0.2) -- (4.5,-0.2) -- (5.7,2.2) -- (0.7,2.2) -- cycle;
    \node at (1.0,1.8) {$(\alpha)$};

    % Phương chiếu Delta
    \draw[thick, blue, ->] (6.0, 3.5) -- (5.0, 1.5) node[right=5pt] {Phương chiếu $\Delta$};

    % Điểm M và tia chiếu
    \fill[red] (2.0, 3.5) circle (2.5pt) node[above right] {$M$};
    \draw[thick, dashed, blue] (2.0, 3.5) -- (1.0, 1.0);
    \fill[blue] (1.0, 1.0) circle (2.5pt) node[below right] {$M'$};

    \node at (2.5, -0.8) {\textbf{Mô hình Phép chiếu song song}};
\end{tikzpicture}
\end{center}

\paragraph{Mục 2: Các tính chất cơ bản của Phép chiếu song song (15 phút)}
\begin{itemize}
    \item \textbf{Tính chất 1 (Bảo toàn tính thẳng hàng):} Phép chiếu song song biến ba điểm thẳng hàng thành ba điểm thẳng hàng và bảo toàn thứ tự giữa các điểm đó.
    $$\text{Nếu } B \text{ nằm giữa } A \text{ và } C \implies B' \text{ nằm giữa } A' \text{ và } C'.$$
    \item \textbf{Tính chất 2 (Biến đường thẳng thành đường thẳng):} Phép chiếu song song biến đường thẳng thành đường thẳng (nếu đường thẳng đó không song song hoặc trùng với phương chiếu), biến tia thành tia, biến đoạn thẳng thành đoạn thẳng.
    \item \textbf{Tính chất 3 (Bảo toàn tính song song):} Phép chiếu song song biến hai đường thẳng song song thành hai đường thẳng song song hoặc trùng nhau.
    $$\text{Nếu } a \parallel b \implies a' \parallel b' \text{ hoặc } a' \equiv b'.$$
    \item \textbf{Tính chất 4 (Bảo toàn tỉ số độ dài):} Phép chiếu song song bảo toàn tỉ số độ dài của hai đoạn thẳng cùng nằm trên một đường thẳng hoặc nằm trên hai đường thẳng song song:
    $$\dfrac{A'B'}{C'D'} = \dfrac{AB}{CD} \quad (AB \parallel CD \text{ hoặc } A, B, C, D \text{ thẳng hàng}).$$
    \textit{Hệ quả đặc biệt:} Phép chiếu song song biến \textbf{trung điểm} của đoạn thẳng thành \textbf{trung điểm} của đoạn thẳng hình chiếu.
\end{itemize}

\begin{scaffoldbox}[Giàn giáo sư phạm cho HS TB - Yếu: Những yếu tố KHÔNG được bảo toàn]
    $\bullet$ \textbf{Cảnh báo sai lầm thường gặp:}\\
    Phép chiếu song song nói chung \textbf{KHÔNG bảo toàn}:\\
    1. Độ lớn của góc (Góc vuông có thể biến thành góc nhọn hoặc góc tù).\\
    2. Độ dài của đoạn thẳng (Đoạn thẳng bị co lại hoặc dãn ra).\\
    3. Tỉ số của hai đoạn thẳng không song song (Đoạn này nghiêng nhiều sẽ co ngắn hơn đoạn kia).\\
    $\bullet$ \textbf{Mẹo ghi nhớ:} Chỉ có \textbf{quan hệ song song}, \textbf{quan hệ thẳng hàng} và \textbf{tỉ số đoạn song song} là được giữ nguyên vẹn!
\end{scaffoldbox}

\subsubsection*{3. Hoạt động 3: Luyện tập (10 phút)}
\begin{itemize}
    \item \textbf{Bài tập 1:} Cho tam giác $ABC$. Gọi $M$ là trung điểm của $BC$. Qua phép chiếu song song lên mặt phẳng $(\alpha)$ theo phương $\Delta$ (không song song với các cạnh tam giác), tam giác $ABC$ biến thành tam giác $A'B'C'$. Điểm $M'$ là hình chiếu của $M$. Hỏi $M'$ có phải là trung điểm của đoạn thẳng $B'C'$ hay không? Vì sao?
    \item \textit{Lời giải:}\\
    Theo Tính chất 4 của phép chiếu song song: Phép chiếu song song bảo toàn tỉ số độ dài của các đoạn thẳng cùng nằm trên một đường thẳng. Vì $M$ là trung điểm của $BC$ nên $\dfrac{MB}{MC} = 1$. Do đó $\dfrac{M'B'}{M'C'} = 1$, suy ra $M'$ chắc chắn là trung điểm của đoạn thẳng $B'C'$.
    \item \textbf{Bài tập 2:} Hình chiếu song song của hai đường thẳng cắt nhau có thể là hai đường thẳng song song được không? Vì sao?
    \item \textit{Lời giải:}\\
    Không thể. Nếu hai đường thẳng $a$ và $b$ cắt nhau tại điểm $I$ thì hình chiếu $a'$ và $b'$ của chúng bắt buộc phải cùng đi qua điểm $I'$ (là hình chiếu của $I$). Do đó $a'$ và $b'$ chỉ có thể cắt nhau tại $I'$ (hoặc trùng nhau nếu mặt phẳng chứa $a, b$ song song với phương chiếu), không bao giờ có thể là hai đường thẳng song song rời nhau.
\end{itemize}

\subsubsection*{4. Hoạt động 4: Vận dụng (05 phút)}
\textbf{Đo chiều cao cây bằng bóng nắng:}\\
Cắm một cây cọc thẳng đứng có chiều dài $1\text{ m}$ trên mặt đất. Đo chiều dài bóng của cọc là $1{,}5\text{ m}$. Cùng thời điểm đó, đo chiều dài bóng của cây cổ thụ trên sân trường là $12\text{ m}$. Học sinh vận dụng tỉ số của phép chiếu song song tính chiều cao của cây cổ thụ:
$$\dfrac{h_{\text{cây}}}{h_{\text{cọc}}} = \dfrac{L_{\text{bóng cây}}}{L_{\text{bóng cọc}}} \implies h_{\text{cây}} = 1 \cdot \dfrac{12}{1{,}5} = 8\text{ m}.$$

\newpage

{\large \textbf{TIẾT 2 (TIẾT 38 THEO PHÂN PHỐI CHƯƠNG TRÌNH)}}\\[0.3em]
\textbf{Nội dung:} Hình biểu diễn của một hình không gian --- Quy tắc vàng vẽ hình --- Năng lực số GeoGebra 3D (Mã 3.1NC1a) --- Năng lực AI Ray-Tracing (Mã 11.C3.2) --- Luyện tập tổng hợp.

\subsubsection*{1. Hoạt động 1: Mở đầu / Khởi động (05 phút)}
\textbf{Chủ đề:} \textit{Vẽ hình hộp chữ nhật trên trang giấy phẳng}
\begin{itemize}
    \item \textbf{a) Mục tiêu:} Giúp học sinh nhận ra nghịch lí: Mặt đáy của hình hộp là hình chữ nhật (các góc đều là $90^\circ$), nhưng trên trang giấy ta lại vẽ nó thành hình bình hành.
    \item \textbf{b) Nội dung:} Tại sao khi vẽ hình hộp chữ nhật hoặc hình lập phương, ta lại vẽ mặt đáy là một hình bình hành chứ không vẽ là hình chữ nhật?
    \item \textbf{c) Sản phẩm:} Học sinh trả lời: Nếu vẽ hình chữ nhật thì góc nhìn trực diện sẽ che khuất các mặt khác, không thấy được chiều sâu ba chiều của vật thể. Phép chiếu song song nghiêng biến hình chữ nhật thành hình bình hành.
    \item \textbf{d) Tổ chức thực hiện:} GV đặt vấn đề $\to$ HS thảo luận $\to$ GV dẫn dắt vào các quy tắc vẽ hình biểu diễn chuẩn mực.
\end{itemize}

\subsubsection*{2. Hoạt động 2: Hình thành kiến thức mới (22 phút)}

\paragraph{Mục 1: Khái niệm Hình biểu diễn của một hình không gian (06 phút)}
\begin{itemize}
    \item \textbf{Định nghĩa:} Hình biểu diễn của một hình $\mathscr{H}$ trong không gian là hình chiếu song song của hình $\mathscr{H}$ trên một mặt phẳng theo một phương chiếu nào đó (hoặc hình đồng dạng với hình chiếu đó).
    \item \textbf{Mục đích:} Biểu diễn chân thực hình dạng, cấu trúc và quan hệ không gian của vật thể 3D lên mặt phẳng 2D (trang giấy, bảng vẽ, màn hình).
\end{itemize}

\paragraph{Mục 2: Quy tắc vàng vẽ hình biểu diễn của các hình phẳng (08 phút)}
Dựa vào các tính chất của phép chiếu song song:
\begin{itemize}
    \item \textbf{Đoạn thẳng, đường thẳng:} Biểu diễn bằng đoạn thẳng, đường thẳng. Giữ nguyên quan hệ song song và tỉ số đoạn thẳng song song.
    \item \textbf{Tam giác:} Hình biểu diễn của một tam giác bất kì (tam giác thường, tam giác cân, tam giác đều, tam giác vuông) đều là \textbf{một tam giác tùy ý}. (Vì phép chiếu không bảo toàn góc và độ dài cạnh).
    \item \textbf{Hình bình hành:} Hình biểu diễn của hình bình hành, hình chữ nhật, hình thoi, hình vuông đều là \textbf{một hình bình hành}.
    \item \textbf{Hình thang:} Hình biểu diễn của hình thang có hai đáy $AB \parallel CD$ là \textbf{một hình thang} có hai đáy tương ứng song song và tỉ số độ dài hai đáy bằng tỉ số của hai đáy ban đầu ($\dfrac{A'B'}{C'D'} = \dfrac{AB}{CD}$).
    \item \textbf{Hình tròn:} Hình biểu diễn của một hình tròn thường là \textbf{một đường elip}.
\end{itemize}

\begin{center}
\begin{tikzpicture}[scale=0.85]
    % Hình biểu diễn hình vuông -> Hình bình hành
    \draw[thick] (0,0) -- (3,0) -- (4,1.5) -- (1,1.5) -- cycle;
    \node at (2,-0.5) {\textbf{Hình bình hành}};
    \node at (2,-1.0) {\small (Biểu diễn cho Vuông, Chữ nhật, Thoi)};

    % Hình biểu diễn đường tròn -> Elip
    \begin{scope}[xshift=7cm]
    \draw[thick] (2, 0.75) ellipse (1.8cm and 0.8cm);
    \fill (2, 0.75) circle (1.5pt) node[below] {$O'$};
    \node at (2,-0.5) {\textbf{Đường Elip}};
    \node at (2,-1.0) {\small (Biểu diễn cho Hình tròn)};
    \end{scope}
\end{tikzpicture}
\end{center}

\paragraph{Mục 3: Tích hợp Năng lực số (Mã 3.1NC1a) (04 phút)}

\begin{pedabox}[Năng lực số (Mã 3.1NC1a): Dựng hình biểu diễn 3D trên GeoGebra 3D]
    $\bullet$ Học sinh mở GeoGebra 3D, chọn công cụ \textbf{Chóp tam giác} (Tetrahedron) và \textbf{Chóp tứ giác}.\\
    $\bullet$ Chuyển đổi chế độ quan sát giữa:
    \begin{itemize}
        \item \textit{Phép chiếu song song trực giao (Orthographic Projection):} Các đường thẳng song song trong không gian luôn song song trên màn hình.
        \item \textit{Phép chiếu phối cảnh (Perspective Projection):} Các đường thẳng song song hội tụ về điểm tụ ở chân trời.
    \end{itemize}
    $\bullet$ Học sinh chụp lại ảnh chụp màn hình hình biểu diễn đúng chuẩn quy ước nét đứt và nét liền.
\end{pedabox}

\paragraph{Mục 4: Tích hợp Năng lực AI (Mã 11.C3.2) (04 phút)}

\begin{aibox}[Tích hợp Năng lực AI (QĐ 2422/QĐ-BGDĐT --- Mã 11.C3.2): AI dò tia (Ray-Tracing) dựng bóng đổ tự động]
    $\bullet$ \textbf{Cơ chế toán học trong đồ họa máy tính và Game 3D:}\\
    Các thuật toán AI xử lí đồ họa hiện đại (NVIDIA DLSS, Unreal Engine 5 Nanite/Lumen) sử dụng kĩ thuật dò tia (Ray-Tracing) dựa chính xác trên \textbf{phép chiếu song song}:\\
    $\bullet$ Ánh sáng mặt trời được mô hình hóa thành chùm tia chiếu song song vô số tia. Thuật toán AI tính giao điểm của từng tia chiếu với bề mặt các hình khối 3D để tự động kết xuất (render) hình chiếu bóng đổ lên mặt đất theo thời gian thực (60 khung hình/giây), mang lại độ chân thực ngoạn mục cho game và phim hoạt hình 3D.
\end{aibox}

\subsubsection*{3. Hoạt động 3: Luyện tập (13 phút)}
\begin{itemize}
    \item \textbf{Bài tập 1 (Thực hành vẽ hình):} Vẽ hình biểu diễn của một hình chóp tứ giác $S.ABCD$ có đáy là hình thang với đáy lớn $AB = 2CD$.
    \item \textit{Hướng dẫn vẽ chi tiết:}\\
    - Bước 1 (Vẽ đáy): Vẽ hình thang $ABCD$ trên mặt giấy sao cho $AB \parallel CD$ và đoạn $AB$ dài gấp đôi đoạn $CD$. Đáy $AB$ thường đặt ở phía sau hoặc phía trước. Đoạn bị khuất vẽ nét đứt (thường là $AD$ và $CD$ tùy góc nhìn).\\
    - Bước 2 (Vẽ đỉnh): Lấy điểm $S$ nằm ngoài mặt phẳng đáy.\\
    - Bước 3 (Nối các cạnh bên): Nối $S$ với các đỉnh $A, B, C, D$. Cạnh nào nhìn thấy vẽ nét liền, cạnh nào bị che khuất vẽ nét đứt (ví dụ $SA$ hoặc $SD$).
    \item \textbf{Bài tập 2 (Nhận biết đúng - sai):} Hình biểu diễn của một hình vuông có thể là một hình vuông được không?
    \item \textit{Lời giải:}\\
    Có thể. Nếu mặt phẳng chứa hình vuông song song với mặt phẳng chiếu thì hình chiếu song song của hình vuông sẽ là một hình vuông bằng chính nó (đây là trường hợp đặc biệt của phép chiếu song song).
\end{itemize}

\subsubsection*{4. Hoạt động 4: Vận dụng (05 phút)}
\textbf{Vẽ phác thảo thiết kế nhà cấp 4 mái Thái:}\\
Vận dụng quy tắc hình biểu diễn (vẽ các cạnh song song bằng các đoạn song song, mặt sàn hình chữ nhật biểu diễn bằng hình bình hành), học sinh phác thảo nhanh khung hình hộp chữ nhật của ngôi nhà kèm theo phần mái dốc chữ A trên trang giấy kẻ ô li.

\newpage

\section*{IV. PHỤ LỤC HỌC LIỆU BỔ TRỢ}

\subsection*{Phụ lục 1: Phiếu học tập số 1 (Tiết 37) --- Tính chất của phép chiếu song song}
\noindent\textbf{Họ và tên:} .................................................................... \textbf{Lớp:} 11A....... \textbf{Nhóm:} ..........

\begin{pedabox}[Nhiệm vụ 1: Hoàn thành bảng tính chất được bảo toàn và không bảo toàn]
Đánh dấu $\checkmark$ vào ô tương ứng:
\begin{center}
\begin{tabularx}{0.95\linewidth}{|p{6.5cm}|X|X|}
\hline
\textbf{Yếu tố hình học} & \textbf{Được bảo toàn} & \textbf{Không bảo toàn} \\
\hline
1. Tính thẳng hàng của ba điểm & & \\
\hline
2. Độ dài đoạn thẳng & & \\
\hline
3. Thứ tự các điểm trên một đường thẳng & & \\
\hline
4. Quan hệ song song của hai đường thẳng & & \\
\hline
5. Độ lớn của góc vuông ($90^\circ$) & & \\
\hline
6. Tỉ số độ dài hai đoạn thẳng cùng nằm trên 1 đường thẳng & & \\
\hline
\end{tabularx}
\end{center}
\end{pedabox}

\subsection*{Phụ lục 2: Phiếu học tập số 2 (Tiết 38) --- Quy tắc vẽ hình biểu diễn}
\noindent\textbf{Họ và tên:} .................................................................... \textbf{Lớp:} 11A....... \textbf{Nhóm:} ..........

\begin{pedabox}[Nhiệm vụ 2: Điền hình biểu diễn phẳng tương ứng]
\begin{center}
\begin{tabularx}{0.95\linewidth}{|p{5.5cm}|X|}
\hline
\textbf{Hình trong không gian thực tế} & \textbf{Hình biểu diễn tương ứng trên mặt phẳng vẽ} \\
\hline
Tam giác vuông, tam giác cân, tam giác đều & Tam giác \dots\dots\dots\dots\dots\dots\dots\dots\dots\dots \\
\hline
Hình chữ nhật, hình vuông, hình thoi & Hình \dots\dots\dots\dots\dots\dots\dots\dots\dots\dots\dots\dots \\
\hline
Hình tròn & Đường \dots\dots\dots\dots\dots\dots\dots\dots\dots\dots\dots\dots \\
\hline
Đoạn thẳng bị che khuất & Vẽ bằng nét \dots\dots\dots\dots\dots\dots\dots\dots\dots\dots \\
\hline
Đoạn thẳng nhìn thấy & Vẽ bằng nét \dots\dots\dots\dots\dots\dots\dots\dots\dots\dots \\
\hline
\end{tabularx}
\end{center}
\end{pedabox}

\subsection*{Phụ lục 3: Bảng Rubric đánh giá năng lực tích hợp trong Bài 14}

\begin{center}
\begin{tabularx}{\linewidth}{|p{3.2cm}|X|X|X|X|}
\hline
\textbf{Tiêu chí} & \textbf{Mức 1 (Chưa đạt)} & \textbf{Mức 2 (Đạt)} & \textbf{Mức 3 (Khá)} & \textbf{Mức 4 (Tốt)} \\
\hline
\textbf{1. Hiểu tính chất phép chiếu song song} & Nhầm lẫn các yếu tố được bảo toàn (tưởng góc vuông được bảo toàn). & Nắm vững các tính chất cơ bản: thẳng hàng, song song, tỉ số đoạn thẳng. & Giải thích thuyết phục vì sao một số yếu tố không bảo toàn. & Vận dụng sắc bén tính chất phép chiếu song song giải các bài toán thực tiễn bóng nắng. \\
\hline
\textbf{2. Kĩ năng vẽ hình biểu diễn không gian} & Vẽ sai quy ước nét đứt nét liền; vẽ hình tròn thành hình tròn méo. & Vẽ đúng hình biểu diễn của tứ diện, hình chóp, hình hộp đơn giản. & Vẽ hình đẹp, chuẩn tỉ lệ, đúng quy tắc elip và hình bình hành. & Biểu diễn các khối hình phức tạp có chiều sâu không gian chân thực, rõ ràng. \\
\hline
\textbf{3. Năng lực số \& AI (Mã 3.1NC1a \& 11.C3.2)} & Chưa biết thao tác phần mềm 3D; chưa liên hệ kiến thức với AI. & Dựng được hình khối cơ bản trên GeoGebra 3D; nêu được ví dụ bóng đổ. & Phân biệt được phép chiếu song song và phối cảnh; hiểu cơ chế AI Ray-Tracing. & Phân tích chuyên sâu vai trò của phép chiếu song song trong đồ họa thời gian thực của Game 3D. \\
\hline
\end{tabularx}
\end{center}

\end{document}
